% FILE: 01_research_question.tex
% PROJECT: experiments-core
% THESIS ROLE: benchmark-experiments-and-equation-verification

\section{Research Question}
\label{sec:experiments-core-rq}

\subsection{Problem Statement}
\label{sec:experiments-core-rq-problem}

The central empirical problem addressed by this project is: \emph{Can the equations
and type laws that govern process conformance be simultaneously verified at three
levels---mathematical (symbolic), executable (Python assertions), and architectural
(Rust compile-time type constraints)---using only artifact evidence that exists in
a local repository, with no mocks, stubs, or simulated results?} A secondary problem
is whether wasm4pm, as an alternative to PM4Py, achieves sufficient performance
and type-safety gains to warrant adoption as the execution engine for board-admissible
process evidence. These problems are not independent: without equation verification,
performance comparisons prove only speed, not correctness; without type-law verification,
performance comparisons prove only that an incorrect computation runs faster.

\subsection{Old-Regime Assumption Challenged}
\label{sec:experiments-core-rq-old-regime}

The old-regime assumption challenged here is that conformance checking is a
\emph{post-hoc analytical activity}: one first runs a process, then one analyses
logs to discover whether it conformed. Under this assumption, conformance is always
retroactive, always approximate, and always disconnected from the runtime type system
of the executing software. PM4Py embodies this assumption: it accepts flat
\texttt{EventLog} objects without requiring that they have passed through any typed
Admission gate, and it returns scalar fitness values without issuing cryptographic
receipts. The old regime treats a fitness score of 0.982 as a number; the new regime
treats it as \texttt{Metric<FITNESS, 982, 1000>} with \texttt{Between01<982, 1000>}
enforced at compile time, signed by Ed25519, and chained to a BLAKE3 model hash.

\subsection{Hypothesis}
\label{sec:experiments-core-rq-hypothesis}

The hypothesis is that the Van der Aalst weighted log fitness equation
$f(L,N) = 1 - \frac{\sum \mathit{freq} \cdot m}{\sum \mathit{freq} \cdot c} - \frac{\sum \mathit{freq} \cdot r}{\sum \mathit{freq} \cdot p}$
and the OCPQ query variable refinement relation
$p \leq_L c \;\triangleq\; \mathrm{dom}(p) \subseteq \mathrm{dom}(c) \;\wedge\; \forall x \in \mathrm{dom}(p),\; p(x) = c(x)$
are jointly necessary and sufficient to characterize lawful conformance in the
object-centric process model, and that wasm4pm's type system, when fully closed
(i.e., when GAP\_001 is resolved), enforces both equations as compile-time invariants
that cannot be bypassed without a type error. The hypothesis further predicts that
wasm4pm achieves at least one order of magnitude performance improvement over PM4Py
in token replay, DFG discovery, and ZKP-compatible conformance tasks.

\subsection{Success Criteria}
\label{sec:experiments-core-rq-success}

Success is defined by three gates, all of which must pass: (1)~all Python validation
scripts execute with zero assertion failures across all test cases---this gate is
currently satisfied as documented in \texttt{checkpoint\_\_experiments\_complete.md};
(2)~the \texttt{audit\_\_experiment\_completeness.md} records 18/18 experiments PASSED
with concrete, non-mock JSON instances---this gate is currently satisfied; and
(3)~GAP\_001 is closed, meaning wasm4pm's conformance API accepts
\texttt{Admission<OcelLog, Ocel20>} rather than flat \texttt{EventLog}, and a
board-admissible \texttt{Receipt} can be manufactured without manual type laundering.
Gate~(3) is not yet satisfied, making the overall status PARTIAL.
